\documentclass{jsarticle}
\usepackage{amsmath}
\usepackage{amssymb}
\begin{document}
\title{オイラーの定数とガンマ関数、\\
そして微分幾何の量子化と虚数による逆三角関数、\\
ベータ関数による場の理論}
\author{Masaaki Yamaguchi}
\date{}
\maketitle

  オイラーの定数は、次の式で成り立っている。
  $$ C = \int{1 \over x^s}dx - \log x $$
  その式を単体積分をすると、
  $$ = \int {\left({\int {1 \over x^s}}dx - \log x\right)}\mathrm{dvol} $$
  ゼータ関数を多重積分すると、　
  $$ = \int \int {1 \over x^s}dx - \int {\log x} \mathrm{dvol} $$
  ここで、対数方程式は、多様体積分がガンマ関数を微分した方程式と
  同値より、
  $$ F = \Gamma = \int e^{-x}x^{1-t}dx $$
  $$ = {d \over df}\int \int {1 \over (y \log y)^{1 \over 2}}dx_m
  - \int e^{-x}x^{1 - t} \log x dx_m $$
  $$ \int x^{1 - t}e^{-x}dV = \int x^{1 - t}dm $$
  $$ \int x^{1 - t}e^{-x}dV = \int x^{1 - t}\mathrm{dvol} $$
  $$ f = \gamma = \Gamma^{'} = \int e^{-x}x^{1 - t} \log x dx $$
  $$ = {d \over d \gamma}\Gamma^{-1} - 
  \int e^{-x}x^{1 - t}\log x dx_m $$
  $$ = {d \over d \gamma}\Gamma^{-1} - \left(\gamma\right)^{\gamma^{'}} $$
  これらより、大域的微分方程式のヒッグス場方程式の逆三角関数の双曲多様体
  に対極する、双曲多様体の余弦定理と同値になる。
  $$ = e^{-f} - e^{f} $$
  $$ = 2 \cos (i x \log x) $$
  結局は、微分幾何の量子化がガンマ関数でもあり、その逆関数はゼータ関数
  でもあり、そして、オイラーの定数を多様体微分をすると、ヒッグス場の
  方程式は正弦定理の逆三角関数であり、このヒッグス場の
  方程式の逆三角関数が、余弦定理の
  双曲多様体として、オイラーの定数になる。
  オイラーの定数は、このガンマ関数から、無理数と証明される。

\end{document}
